\section{Forecasting methodology}\label{sec:method}

\paragraph{Features.}
Each row is a report of entity $e$ at origin $t$ and carries 24 base features plus one horizon-specific feature, all computed from reports dated $\le t$ (Table~\ref{tab:features}).
Missing values are set to zero.
Every feature that can be missing has a companion indicator or count (\texttt{has\_ly}, \texttt{cnt30}), so no fitted imputation statistic is needed and none can leak.

\begin{table}[t]
\centering\small
\caption{Feature groups ($\ell=\ln P$). All frames end at the origin.}\label{tab:features}
\begin{tabular}{@{}lp{0.74\linewidth}@{}}
\toprule
Group & Features \\
\midrule
Level / lags & $\ell_t$; $\ell_t-\ell_{t^{(-k)}}$ for $k\in\{1,2,5\}$ previous reports; days since previous report \\
Rolling (7/30/90 d) & deviation from the rolling mean; std (30, 90); report counts (30, 90); distance to 30-day max/min \\
Range & $(\max-\min)/\text{modal}$ on day $t$ \\
Seasonal & deviation from the level at $t-365\pm3$; indicator; last year's change from $t-365$ to $t+h-365$ \\
Cross-section & same-day mean 30-day deviation and 1-report return of the commodity nationally and within the state; number of reporting markets \\
Calendar & day-of-year (sine, cosine); day of week \\
Categorical & commodity index (fitted on training rows) \\
\bottomrule
\end{tabular}
\end{table}

\paragraph{Baselines.}
Five baselines are reported, all as log returns:
\begin{itemize}
  \item persistence ($\hat y=0$);
  \item geometric moving averages over 7 and 30 days;
  \item seasonal naive (last year's level around $t+h-365$; persistence when unavailable);
  \item seasonal-adjusted persistence (today's level times last year's change over the same horizon).
\end{itemize}

\paragraph{Learners.}
All learners are Spark MLlib models, one per horizon (direct strategy):
\begin{itemize}
  \item ridge regression with squared and Huber losses;
  \item a random forest (60 trees, depth 10);
  \item gradient-boosted trees with squared loss (depth 6, 200 iterations, step 0.05);
  \item gradient-boosted trees with absolute loss in three configurations.
\end{itemize}
MLlib's absolute-loss boosting fits each tree after the first to $\operatorname{sign}$ of the residual, without a leaf line search.
Each iteration therefore moves the prediction by up to the step size, which is large relative to typical log returns of about 0.01.
We therefore train these models on the target in percent with step sizes of 0.25--0.5, and rescale inside the persisted pipeline.
A unit test confirms that this pipeline recovers the conditional median of a zero-inflated target, whereas squared loss recovers the mean.

\section{Experimental protocol}\label{sec:protocol}

\paragraph{Periods.}
\begin{itemize}
  \item Training origins: 2016-01-01 to 2023-12-31.
  \item Validation origins: 2024.
  \item Test origins: 2025-01-01 to 2026-04-21.
\end{itemize}
A row is kept in a period only if its target date also lies in that period.
This purge removes the one path by which a label can cross a boundary, and a unit test checks it.
Training uses a uniform random sample of $2\times10^6$ rows per horizon (seed fixed), drawn from 15--16M eligible rows.
Hyper-parameters are tuned on a random sample of $6\times10^5$ validation rows.

\paragraph{Selection and amendments.}
For each horizon, the family with the lowest validation MAE (in rupees) is selected, baselines included.
The search space was amended twice, before any test metric was read, on validation evidence only:
\begin{itemize}
  \item first, because squared-loss learners tied persistence at the next-report horizon, where 34.4\% of target changes are exactly zero;
  \item second, because MLlib's absolute loss is mis-scaled for log returns, as described above.
\end{itemize}
The logs of both aborted runs are kept in the repository.

\paragraph{Metrics.}
The following are computed on the rupee scale over all test rows:
\begin{itemize}
  \item MAE, RMSE, WAPE $=\sum|e|/\sum|A|$, sMAPE, and bias (mean signed error);
  \item MASE, scaled by each entity's in-sample $h$-step naive MAE on training rows \cite{hyndman2006another};
  \item relative MAE against persistence on identical rows;
  \item directional accuracy on rows with $|y|\ge1\%$.
\end{itemize}
MAPE is not used: low-priced leafy vegetables make its denominators unstable.
Diebold--Mariano tests \cite{diebold1995comparing} compare each model with persistence on the series of daily mean absolute-error differences, using a Newey--West variance (lag $=h$) and the Harvey--Leybourne--Newbold correction.
Errors are aggregated per date first because same-day errors are strongly correlated across markets.
We also report the share of series (those with at least 30 test rows) on which the selected model beats persistence.

\paragraph{Intervals.}\label{sec:intervals}
For the selected model we use split-conformal intervals with Mondrian bins: validation residuals $r=y-\hat y$ are grouped into five volatility bins by 30-day log-price standard deviation, with bin edges taken from validation quantiles.
The 5/10/90/95\% residual quantiles of each bin give 80\% and 90\% intervals.
Exchangeability does not hold for time series, so coverage is measured on the test period rather than assumed \cite{gibbs2021adaptive,xu2023conformal}.

\paragraph{Stability.}
A rolling-origin backtest refits the selected configuration on expanding windows ending in 2020--2023 and tests on each following year.

\paragraph{Anomaly detection.}
No labels exist for real anomalies.
Instead, corruptions are injected into 1\% of \num{400000} sampled 2025 test rows, half of each kind:
\begin{itemize}
  \item unit or keying errors ($\times10$, $\times100$, $\div10$);
  \item shocks of $\pm35$--60\%.
\end{itemize}
We compare the production rule (trailing robust $z$ beyond 3.5 with a move of more than 20\%, or more than a factor of five) with two baselines: a previous-report change and the legacy commodity-IQR rule, refitted on training data.
